% Generated draft: requires mathematical and source review.
\documentclass[11pt]{article}
\usepackage[margin=25mm]{geometry}
\usepackage{fontspec}
\setmainfont{Latin Modern Roman}
\usepackage{amsmath,amssymb,mathtools,unicode-math}
\setmathfont{Latin Modern Math}
\usepackage{enumitem,needspace,xurl,hyperref}
\setlength{\parindent}{0pt}
\setlength{\parskip}{5pt}
\setlength{\emergencystretch}{3em}
\newcommand{\N}{\mathbb{N}}
\newcommand{\Z}{\mathbb{Z}}
\newcommand{\Q}{\mathbb{Q}}
\newcommand{\R}{\mathbb{R}}
\newcommand{\C}{\mathbb{C}}
\newcommand{\F}{\mathbb{F}}
\newcommand{\A}{\mathbb{A}}
\newcommand{\PP}{\mathbb{P}}
\newcommand{\E}{\mathbb{E}}
\newcommand{\eps}{\varepsilon}
\newcommand{\dd}{\,\mathrm{d}}
\begin{document}
\section*{UnsolvedMath ID: 30002877}
\textbf{Source problem code:} catalogue entry
\subsection{Definitions and mathematical statement}
Let \((X,\mathcal O_X)\) be a connected ringed space and let \(A\) be an Azumaya \(\mathcal O_X\)-algebra, as in the supplied statement. The source does not specify the class of ringed spaces or the precise convention for Azumaya algebras. The following gives the usual interpretation with a commutative structure sheaf and well-defined positive locally constant ranks; its applicability to the intended setting requires confirmation.

Under this interpretation, an Azumaya algebra is a sheaf of associative unital \(\mathcal O_X\)-algebras that is locally free of finite positive rank as an \(\mathcal O_X\)-module and for which the canonical map
\[
 A\otimes_{\mathcal O_X}A^{\mathrm{op}}
 \longrightarrow \mathcal E nd_{\mathcal O_X}(A),
 \qquad a\otimes b\longmapsto (c\longmapsto acb),
\]
is an isomorphism. Here \(A^{\mathrm{op}}\) denotes the opposite algebra. Its degree is the positive integer \(n\) such that
\[
 \operatorname{rank}_{\mathcal O_X}(A)=n^2.
\]
Connectedness makes this degree constant in the usual nonempty setting. The source does not explicitly exclude the empty space.

Write \([A]\) for the Brauer class of \(A\). Under the usual stabilization convention, two Azumaya algebras \(A\) and \(A'\) are Brauer-equivalent if there exist finite locally free \(\mathcal O_X\)-modules \(E,F\) of positive rank and an isomorphism of \(\mathcal O_X\)-algebras
\[
 A\otimes_{\mathcal O_X}\mathcal E nd_{\mathcal O_X}(E)
 \cong
 A'\otimes_{\mathcal O_X}\mathcal E nd_{\mathcal O_X}(F).
\]
Define the degree spectrum of the class by
\[
 D_X([A])=
 \left\{n\in\Z_{>0}:\begin{array}{l}
 \text{there exists an Azumaya algebra }A'\text{ on }X,\\
 [A']=[A]\text{ and }\operatorname{rank}_{\mathcal O_X}(A')=n^2
 \end{array}\right\}.
\]
This set is nonempty because it contains \(\deg(A)\). In the convention of the supplied record, its index is
\[
 \operatorname{ind}([A])=\gcd D_X([A]),
 \qquad
 \operatorname{ind}([A])\Z
 =\left\langle n:n\in D_X([A])\right\rangle_{\Z}.
\]

The problem is to determine \(D_X([A])\) for an arbitrary connected ringed space \(X\) and Azumaya algebra \(A\) in the intended category, and to characterize its greatest common divisor in terms of \(X\) and \([A]\). No particular form of such a characterization is specified. The definition by a greatest common divisor does not assert that the index itself occurs as the degree of a representative. The source leaves the category of spaces and any additional hypotheses unresolved; no restriction to schemes, topological spaces with a particular structure sheaf, or fields is supplied.
\subsection{Short English statement}
For a fixed Azumaya algebra on a connected ringed space, determine exactly which positive integers occur as degrees of Brauer-equivalent Azumaya algebras, and characterize their greatest common divisor, called the index in the supplied record. The intended category of ringed spaces and the desired form of the characterization need clarification.
\subsection{Sources}
\begin{itemize}
\item Local UnsolvedMath registry, permanent ID \(30002877\), title: \emph{Degree Spectra and Index of Brauer-Equivalent Azumaya Algebras}; exact source problem code: \texttt{catalogue entry}. The record belongs to the Oberwolfach Reports Open Problems collection. No individual proposer is supplied; the recorded proposed year is 2015.
\item Supplied report citation: \emph{Cohomology of Finite Groups: Interactions and Applications} (2016), Oberwolfach Reports, pp.~1317--1384, DOI \texttt{10.4171/owr/2015/24}; \url{https://doi.org/10.4171/owr/2015/24}. The supplied background describes these reports as documenting mathematical workshops and their open-problem discussions. No precise location of this problem within the cited page range is supplied.
\item Supplied literature reference: Perry, \emph{Period-index bounds for Brauer classes}; \url{https://arxiv.org/abs/2608.03684}.
\item Supplied literature reference: Antieau--Williams, \emph{The period-index problem for twisted topological K-theory}; \url{https://arxiv.org/abs/2003.10539}.
\item Literature assessment dated exactly \texttt{2026-08-21}: the supplied curation labels the problem open, reports that period-index methods give important bounds and examples while the general degree question remains unresolved, and states that a conjectural universal bound is now known to fail. This is supplied dated curation derived from the OWR record, not an independently verified literature assessment. The bound and its hypotheses are not specified. Time-sensitive status and cited claims require independent verification.
\item The registry labels the statement \texttt{corrected\_verified} and says that its mathematical content was preserved while notation, typography, or omitted definitions were repaired for publication. These are supplied editorial assertions, not a mathematical verification of this draft.
\end{itemize}
\end{document}
